\section{消融实验结果}

为验证 AGF-ADNet 中各组成模块的有效性，本文在相同训练集与测试集上设置消融实验。训练数据为正常工况样本，测试数据为 C5 故障样本；检测阈值均由训练集异常分数确定，未使用测试标签参与阈值选择。评价指标包括 Accuracy、Precision、Recall、FDR、FRA 和 F1-Score。其中，FDR 按故障检出率计算，与 Recall 均为 \(TP/(TP+FN)\)；FRA 为误报率，即 \(FP/(FP+TN)\)。完整模型结果采用前文 AGF-ADNet 分类实验结果，消融模型结果来自 \texttt{test/outputs/} 下各实验输出。

\begin{table}[!htbp]
    \centering
    \caption{AGF-ADNet 消融实验结果}
    \label{tab:agf_adnet_ablation_results}
    \begin{tabular}{p{4.3cm} p{1.6cm} p{1.6cm} p{1.6cm} p{1.4cm} p{1.4cm} p{1.7cm}}
        \toprule
        \textbf{模型设置} & \textbf{Accuracy} & \textbf{Precision} & \textbf{Recall} & \textbf{FDR} & \textbf{FRA} & \textbf{F1-Score} \\
        \midrule
        AGF-ADNet（完整模型） & 0.9910 & 1.0000 & 0.9820 & 0.9820 & 0.0000 & 0.9910 \\
        图频门控 + Transformer & 0.9303 & 1.0000 & 0.8605 & 0.8605 & 0.0000 & 0.9250 \\
        图频门控 + 直接检测 & 0.6330 & 1.0000 & 0.2660 & 0.2660 & 0.0000 & 0.4202 \\
        仅图专家 + Transformer & 0.8673 & 0.9939 & 0.7390 & 0.7390 & 0.0045 & 0.8477 \\
        仅频域专家 + Transformer & 0.5363 & 0.9677 & 0.0750 & 0.0750 & 0.0025 & 0.1392 \\
        仅图专家 + 直接检测 & 0.5000 & 0.0000 & 0.0000 & 0.0000 & 0.0000 & 0.0000 \\
        仅频域专家 + 直接检测 & 0.5000 & 0.0000 & 0.0000 & 0.0000 & 0.0000 & 0.0000 \\
        \bottomrule
    \end{tabular}
\end{table}

从表~\ref{tab:agf_adnet_ablation_results} 可以看出，完整 AGF-ADNet 在所有模型中取得最优结果，Accuracy 为 0.9910，F1-Score 为 0.9910，Precision 为 1.0000，FRA 为 0.0000。这说明完整模型在该测试集上没有产生正常样本误报，同时能够检出绝大多数故障窗口。与最接近完整模型的“图频门控 + 普通 Transformer”相比，完整模型的 Accuracy 提高 0.0607，Recall 和 FDR 提高 0.1215，F1-Score 提高约 0.0660。该差异表明，完整模型中的采样率感知检测器、相位信息建模以及分布漂移分数对多采样率过程的故障识别具有明显贡献。普通 Transformer 虽然能够利用图频融合后的完整序列进行时序重构，但没有显式利用不同变量的采样节拍，因此对部分故障窗口的响应仍不充分。

图频门控融合对检测性能有直接提升。当保留图专家、频域专家和门控融合，但去除 Transformer 检测器，仅使用融合重构误差进行直接检测时，模型 Accuracy 为 0.6330，Recall 为 0.2660，F1-Score 为 0.4202。该结果高于两个单专家直接检测模型，说明图结构信息与频域信息的互补能够增强异常分数的区分度。然而，该模型仍有大量故障窗口未被检出，表明仅依赖插补重构误差不足以稳定刻画故障在时间维度上的连续演化。换言之，门控插补能够改善数据补全质量，但异常检测仍需要专门的时序判别模块。

单专家实验进一步说明了图结构专家和频域专家的作用差异。在加入 Transformer 检测器的情况下，仅图专家模型取得 0.8673 的 Accuracy 和 0.8477 的 F1-Score，明显优于仅频域专家模型的 0.5363 Accuracy 和 0.1392 F1-Score。这说明对于 C5 故障，变量间相关关系和空间依赖结构是更主要的判别信息来源。仅频域专家模型的 Precision 仍达到 0.9677，但 Recall 只有 0.0750，说明它只对少量异常窗口给出高置信度告警，漏检较多。相比之下，图专家能够提供更连续的故障响应，但仍存在少量误报，FRA 为 0.0045。

双专家门控与 Transformer 检测器结合后，模型性能进一步提升。“图频门控 + 普通 Transformer”的 Accuracy 达到 0.9303，Recall 达到 0.8605，F1-Score 达到 0.9250，并且 FRA 降为 0.0000。相较于“仅图专家 + Transformer”，其 Recall 提高 0.1215，F1-Score 提高约 0.0773；相较于“仅频域专家 + Transformer”，提升更加显著。这表明频域专家虽然单独使用时检出能力有限，但与图专家通过门控机制融合后，可以补充图结构建模难以覆盖的周期性和局部频谱变化信息。门控权重统计也支持这一点：在测试集上，“图频门控 + 普通 Transformer”的频域专家平均权重约为 0.7364，图专家平均权重约为 0.2636，且低采样率变量上的频域权重更高。这说明门控层并非简单平均两个专家输出，而是根据变量采样特征和重构状态动态分配两类专家的贡献。

直接检测类模型的结果说明，后处理可以降低误报，但不能弥补异常分数本身区分能力不足的问题。仅图专家和仅频域专家的直接检测模型最终均未形成有效故障告警，Accuracy 均为 0.5000，Recall 和 F1-Score 均为 0。由于测试集正常与异常窗口数量相同，该结果等价于模型几乎将全部样本判为正常。结合输出文件可知，两个模型的原始异常分数中存在少量越过阈值的点，但这些点持续时间较短，经最短异常持续长度约束后被抑制。这一现象说明，单一插补专家产生的重构误差容易出现零散波动，难以形成稳定、连续的故障片段。

综上，消融实验表明 AGF-ADNet 的性能提升来自多个模块的协同作用：图专家负责刻画变量间依赖关系，频域专家补充多采样率序列中的周期性和频谱信息，门控机制实现两类专家的自适应融合，Transformer 检测器进一步建模故障的时序演化；在此基础上，采样率感知特征、相位信息和分布漂移分数显著增强了完整模型对故障窗口的覆盖能力。各模块缺失时，模型主要表现为漏检增加，而完整模型在保持零误报的同时取得最高检出率，说明所提出结构对于多采样率过程故障检测是有效且必要的。
